% Long justified lines in Type 1 fonts, as in real papers: T1 Latin Modern,
% whose widths are no whole thousandths of an em (e is 444.45), roman, italic,
% bold and small, and the same text in Latin Modern's OpenType fonts. Each
% glyph is where TeX put it, to the end of every line.
\documentclass{article}
\usepackage[T1,TU]{fontenc}
\usepackage{lmodern}
\begin{document}
\fontencoding{T1}\selectfont
Vlasov equations arise in various applications, including the social sciences,
biology, and plasma physics; specifically, the equation appears in models for
muscle contraction and finance, and in particle-in-cell methods for the
numerical solution of kinetic equations. We note that our results extend to
time-dependent coefficients without much change, and that the constants in the
estimates depend only on the Lipschitz constants of the coefficients.

\textit{In italic: the estimates of the previous section are uniform in the
mesh size, so the scheme converges at the rate the regularity of the solution
allows, and the constants depend only on the bounds assumed for the
coefficients and on the final time.} \textbf{In bold: the main theorem holds
for every finite time interval, and its proof uses nothing beyond Gronwall's
inequality and a careful count of the terms in the expansion.}

{\small A smaller size has other widths again: the harmonic series diverges,
though very slowly, its partial sums growing like the logarithm of the number
of terms, while its terms tend to zero, which is not enough for convergence.\par}

\fontencoding{TU}\selectfont
The same words in OpenType: Vlasov equations arise in various applications,
including the social sciences, biology, and plasma physics; specifically, the
equation appears in models for muscle contraction and finance, and in
particle-in-cell methods for the numerical solution of kinetic equations.
\end{document}
